\chapter{从零手写自动微分与反向传播引擎 (Micrograd)}

\section{Lab 03：从零手写自动求导（Autograd）}

\begin{noteBox}[核心导读与背景]
对应理论：\textbf{第 3 篇：神经网络如何学习：导数、链式法则、反向传播与交叉熵}
前置：读完第 3 篇第一 \textasciitilde{} 六节。
依赖：纯 Python 标准库。
\end{noteBox}

\section{为什么要自己写一遍？}

PyTorch 的 \texttt{loss.backward()} 是整个深度学习最大的"黑盒"。只要你亲手写过一次，就会明白：
\begin{itemize}
  \item 反向传播 = \textbf{链式法则 + 拓扑排序 + 梯度累加}，没有任何魔法；
  \item 为什么每一步训练前必须 \texttt{zero\_grad()}（梯度是 \texttt{+=} 累加的）；
  \item 为什么训练要保存激活值（\texttt{\_backward} 闭包里用到了前向时的 \texttt{self.data}）；
  \item 为什么推理时可以用 \texttt{torch.no\_grad()} / \texttt{torch.inference\_mode()} 省掉大量显存（不建图、不存激活）。
\end{itemize}

\section{运行}

\begin{lstlisting}[language=bash]
python3 lab03_autograd_from_scratch/micrograd_from_scratch.py
\end{lstlisting}

\begin{center}
\small
\begin{tabularx}{\textwidth}{p{3cm} p{4.5cm} X}
\toprule
\textbf{实验} & \textbf{验证的结论} & \textbf{对应章节} \\
\midrule
1 & 手算表格 dL/dw = −12, dL/db = −6 & 第 3 篇 §2.3 \\
2 & 自动求导 = 数值差分（梯度检验） & 第 3 篇 §3 \\
3 & Softmax 雅可比 $p_i(\delta_{ij}-p_j)$；交叉熵梯度 $p - y$；logits 放大后梯度坍缩 & 第 3 篇 §6、第 4 篇 §3 \\
4 & $\partial L/\partial W = X^T G$，$\partial L/\partial X = G W^T$ & 第 3 篇 §4 \\
5 & Sigmoid 连乘梯度消失 vs 残差畅通 & 第 3 篇 §7、第 4 篇 §5 \\
6 & 单层学不会 XOR，两层可以 & 第 1 篇 §4.2、第 0 篇族谱 \\
\bottomrule
\end{tabularx}
\end{center}

\section{动手练习（做完才算过关）}

\begin{enumerate}
  \item 给 \texttt{Value} 加一个 \texttt{silu()} 算子（$x\sigma(x)$），自己推导导数，并用实验 2 的方法做梯度检验。
  \item 把实验 6 的激活函数从 \texttt{tanh} 换成 \texttt{sigmoid}，学习率不变，观察收敛速度变化，并用第 3 篇 §7 的导数上界解释。
  \item 把实验 5 的深度改成 30，看纯 Sigmoid 的梯度变成多少。
  \item （进阶）跟着 \href{https://karpathy.ai/zero-to-hero.html}{Karpathy 的 micrograd 视频} 实现 \texttt{Neuron / Layer / MLP} 类，并训练一个二分类的"月牙"数据集。
\end{enumerate}

\section{学完之后}

你已经可以放心使用 \texttt{torch.autograd} 了——它做的事和这 350 行完全一样，只是算子换成了张量、实现换成了 C++/CUDA。
下一步：\textbf{第 5 篇} $\rightarrow$ \textbf{Lab 05：亲手训练一个迷你 GPT}。


\section{实验核心源码精选与解析}

本实验配套有完整可运行的工业级验证代码，重点实现细节如下：


\subsection{源码实现：\texttt{micrograd\_from\_scratch.py}}

\lstinputlisting[language=Python, caption={\texttt{micrograd\_from\_scratch.py}}]{code/lab03_autograd_from_scratch_micrograd_from_scratch.py}

